\documentclass[12pt]{article}

\usepackage[spanish]{babel}
\usepackage[utf8]{inputenc}
\usepackage{amsmath, amssymb}
\usepackage{geometry}
\usepackage{graphicx}
\usepackage{booktabs}

\geometry{margin=1in}

\title{Módulo 2: Identificación y Medición de Riesgos Previsionales}
\author{MACROPOL}
\date{}

\begin{document}

\maketitle

\section{Contexto}

El stress testing en sistemas de pensiones de capitalización individual requiere una adecuada identificación y medición de los riesgos previsionales. A diferencia del sistema bancario, donde el foco es la solvencia, en pensiones el análisis se centra en la suficiencia de las pensiones futuras y en la estabilidad del ahorro acumulado.

Los sistemas previsionales están expuestos a riesgos múltiples y correlacionados, cuya materialización puede comprometer la sostenibilidad del sistema y el bienestar de los afiliados.

\textbf{Idea central:} el stress testing es tan robusto como la medición de los riesgos que incorpora.

\section{Modelo / Herramienta}

\subsection{Clasificación de riesgos}

\begin{itemize}
\item \textbf{Riesgo de mercado:} variaciones en precios de activos financieros.
\item \textbf{Riesgo de longevidad:} aumento en la esperanza de vida.
\item \textbf{Riesgo macroeconómico:} cambios en crecimiento, inflación y tasas de interés.
\item \textbf{Riesgo de liquidez:} necesidad de liquidar activos en condiciones adversas.
\end{itemize}

\subsection{Dinámica del fondo de pensiones}

\begin{equation}
W_{t+1} = W_t (1 + r_t) + c_t
\end{equation}

donde:
\begin{itemize}
\item $W_t$: saldo acumulado
\item $r_t$: retorno del portafolio
\item $c_t$: contribuciones
\end{itemize}

\subsection{Tasa de reemplazo}

\begin{equation}
RR = \frac{\text{Pensión}}{\text{Último salario}}
\end{equation}

\subsection{Modelación de riesgos}

\subsubsection{Riesgo de mercado}

\begin{equation}
r_t = \ln\left(\frac{P_t}{P_{t-1}}\right)
\end{equation}

Modelo de volatilidad (GARCH):

\begin{equation}
\sigma_t^2 = \alpha_0 + \alpha_1 \epsilon_{t-1}^2 + \beta \sigma_{t-1}^2
\end{equation}

\subsubsection{Riesgo macroeconómico}

\begin{equation}
Y_t = A Y_{t-1} + \epsilon_t
\end{equation}

donde $Y_t$ incluye PIB, inflación y tasas de interés.

\subsubsection{Riesgo de longevidad}

Modelo Lee-Carter:

\begin{equation}
\ln(m_{x,t}) = a_x + b_x k_t + \epsilon_{x,t}
\end{equation}

\subsection{Interacción de riesgos}

\begin{equation}
\Sigma =
\begin{bmatrix}
\sigma^2_{r} & \rho_{r,y} \\
\rho_{r,y} & \sigma^2_{y}
\end{bmatrix}
\end{equation}

Los riesgos están correlacionados, lo que amplifica los efectos en escenarios adversos.

\section{Implementación}

\subsection{Simulación de retornos}

\begin{verbatim}
returns <- diff(log(prices))
volatility <- sd(returns)
\end{verbatim}

\subsection{Modelo VAR}

\begin{verbatim}
library(vars)
model <- VAR(data, p=2)
forecast <- predict(model, n.ahead=12)
\end{verbatim}

\subsection{Simulación Monte Carlo}

\begin{verbatim}
sim_returns <- rnorm(1000, mean=mean(returns), sd=sd(returns))
wealth <- cumprod(1 + sim_returns)
\end{verbatim}

\section{Ejercicio Aplicado}

\textbf{Escenario adverso:}
\begin{itemize}
\item Caída de mercado: -20\%
\item Aumento de tasas: +200 pb
\item Incremento en longevidad: +2 años
\end{itemize}

\textbf{Tareas:}
\begin{itemize}
\item Simular la trayectoria del fondo
\item Calcular la tasa de reemplazo
\item Comparar con escenario baseline
\end{itemize}

\section{Interpretación}

\subsection{Baseline vs adverso}

\begin{center}
\begin{tabular}{lcc}
\toprule
Variable & Baseline & Adverso \\
\midrule
Retornos & Normales & Negativos \\
Longevidad & Estable & Aumenta \\
PIB & Crecimiento & Recesión \\
\bottomrule
\end{tabular}
\end{center}

\subsection{Insights}

\begin{itemize}
\item El riesgo relevante es conjunto, no individual
\item La tasa de reemplazo es altamente sensible a shocks
\item La correlación entre riesgos amplifica pérdidas
\end{itemize}

\section{Implicaciones Regulatorias}

El análisis de riesgos previsionales es fundamental para:

\begin{itemize}
\item Supervisión de fondos de pensiones
\item Diseño de portafolios de inversión
\item Reformas previsionales
\end{itemize}

El regulador no solo evalúa pérdidas financieras, sino el impacto sobre el bienestar previsional.

\section{Mensaje Final}

\begin{center}
\textit{El stress testing no comienza con escenarios, sino con una comprensión profunda de los riesgos.}
\end{center}

\end{document}